\section{Three-twist refinement and tight purity extraction}
\label{sec:resolved}
\begin{theorem}\label{thm:resolved}
Assume the pinned spatial-transfer and sector identities of
\cite{upstream}, its recorded untwisted and $P$-twisted thermal integer
totals and 120-site trial, the retained 88-site trial, and the new
cyclic-twist integer totals certified below. At coupling $J=1$,
\[
 \boxed{\gap_L>\frac{29}{200}=0.145
 \quad\text{for every integer }L\ge24.}
\]
The ground state is unique. The strict bounds $0.02$, $0.09$, and
$0.12$ hold for every integer length at least $18$, $20$, and $22$,
respectively. Separately, $\gap_{18}>0.08$.
Thus $\liminf_{L\to\infty}\gap_L\ge0.145$, with no parity restriction.
\end{theorem}

This revision supplies the complete cyclic-twist table at lengths
$4,\ldots,12$, separates $I$ and $O$ already at the input temperature,
and uses a local energy inequality to reduce Horner rounding errors.
It retains the sector recurrence proved in the next section. A sharper
elementary conversion of purity into a gap improves the finite-length
coverage. These are quantitative refinements conditional on the cited
analytic premises; neither priority nor a new general moment method is
claimed. Earlier statements and certificates are retained unchanged.

\subsection{A local energy bound and geometric rounding errors}
\begin{lemma}\label{lem:local-energy}
For the five-site open spin-one Hamiltonian $H^{\mathrm{open}}_5$,
\[
 H^{\mathrm{open}}_5>-\frac{35}{6}I.
\]
Consequently $H_L(g)\ge-35L/24$ for $L\ge5$ and each of the three
boundary rotations $g=1,P,C$ used here.
\end{lemma}
\begin{proof}
In the spin-weight basis, the open Hamiltonian has diagonal
$\sum_{j=1}^{4}w_jw_{j+1}$ and unit entries for allowed opposite unit
steps on adjacent spins. Its nonnegative-magnetization block dimensions
are $51,45,30,15,5,1$. Spin reversal supplies the remaining blocks.
Fraction-free Bareiss elimination verifies that every leading principal
minor of the integer matrix $6H^{\mathrm{open}}_5+35I$ is positive.
The checker verifies every division exactly; its complete pivot lists
are in \path{independent_results/local_energy.json}.
Sylvester's criterion gives strict positivity in each block.

Sum the translated five-site open blocks around a ring. Each bond is
counted four times. For a twisted ring, any such open path is unitarily
equivalent to the untwisted path: on-site rotations move its single
seam to the missing bond. Thus the same local inequality applies, and
division by four gives the asserted ring bound.
\end{proof}

The thermal calculation retains the upstream matrix
$R=I-(H_n(g)+3nI/2)/16$. For $5\le n\le12$, the lemma and
$H_n(g)\le nI$ imply
\[
 \operatorname{spec}(R)\subseteq
 [1-5n/32,\,1-n/384],\qquad
 \|R\|\le\kappa_n=1-n/384<1.
\]
At $n=4$ we keep the original bound $\|R\|\le1$.
This stronger contraction bound changes only the error estimate, not
the recorded integer columns.

For completeness, set $Q_*=2^{55}$, $J_*=280$, and
$p_j=98^j/j!\big/\sum_{k=0}^{280}98^k/k!$.
The exact Horner coefficients are $\lfloor Q_*p_j\rfloor$; only indices
through 191 are nonzero. The conditioned Poisson-tail error is less
than $3^{98}/2^{280}<Q_*^{-1}$.
For a magnetization block of dimension $d$, put
$d^+=\lfloor\sqrt d\rfloor+1$.
A real-coordinate floor has Euclidean error at most $\sqrt d$.
For the cyclic twist, use $u+v\omega$, with
$\omega=e^{2\pi i/3}$; its squared modulus is $u^2-uv+v^2$.
Both floor errors lie in $[-1,0]$, where
$a^2-ab+b^2\le1$, so the same vector error bound applies.
One coefficient floor contributes at most one more unit.
Geometric propagation therefore gives the integer column-error bounds
\[
 e_{n,d}=\begin{cases}
 281(d^++1)+1,&n=4,\\
 \left\lceil384(d^++1)/n\right\rceil+1,&5\le n\le12.
 \end{cases}
\]
The final unit includes the Poisson approximation. If $d_M$ is the
dimension of magnetization block $M$, use
\[
 E_n=\left\lfloor\sqrt{\sum_{M=-n}^{n}d_M e_{n,d_M}^2}\right\rfloor+1
\]
as the Frobenius error in $Q_*$ units. Repeating representative columns
with their orbit multiplicities preserves this estimate: equality of
the \emph{exact} column norms follows from the upstream symmetries,
whereas the rounded columns are bounded individually.
The resulting trace endpoints have the same form as in the retained
derivation, with this smaller $E_n$.

\subsection{Extended cyclic table and resolved input sectors}
The new evaluator \path{compute_thermal_moments.cpp} implements the
same integer Horner construction for $g=1,P,C$. It uses denominator
32 also at odd lengths; the cyclic phase is represented by two integer
coordinates. Ordinary column operations fit signed 64-bit integers:
the absolute coordinate row sums are at most 80 and
$80(2Q_*+2\max_d e_{n,d})<2^{63}$.
The per-column quadratic sums and their orbit-weighted values fit
signed 128-bit integers. A checked four-limb accumulator stores the
total, whose conservative upper bound is below $2^{256}$.
The coefficient files are generated as exact rational floors.

All nine cyclic totals are recorded in
\path{independent_results/extended_cyclic_totals.json}.
They were recomputed from the distributed C++ source. At lengths
4, 6, and 8, they agree exactly with the unmodified upstream packed
engine using the same coefficients. The untwisted and $P$-twisted
four-site totals were cross-checked too. At odd length five, an
independent all-column calculation without orbit reduction gives
overlapping rigorous norm intervals. These checks support the
implementation; they are not independent peer review of the proof.

Write $C_n=Z_n(b_*,C)$ and $b_*=49/4$. The sector inversion is
\[
 I_n=\frac{J_n+2C_n}{3},\qquad
 O_n=\frac{J_n-C_n}{3},\qquad
 Z_n=I_n+2O_n+3N_n.
\]
The old untwisted and $P$-twisted totals are re-enclosed using the same
geometric error bound. Fixed quartic filters $x^4P(x)^2$ then give
\begin{align*}
 -.59564&<\mu_I<1.0008406,\\
 -.577508&<\mu_O<.699301,\\
 -.8710233&<\mu_N<.7502156.
\end{align*}
The coefficient vectors are fixed in
\path{verify_resolved_refinement.py}. Every coefficient after shifting
the relevant exterior ray is positive, and the endpoint filter exceeds
its rigorously bounded moment sum. Thus these are full-ray exclusions.

For each sector the positive measure
$\sum_j\mu_j^4\delta_{\mu_j}$ permits joint use of moments 4 through 12.
Fixed rational majorants and minorants in
\path{independent_results/resolved_moment_polynomials.json}
are certified on the entire respective support by Bernstein subdivision.
Their signed odd-moment bounds are used only for signed traces;
odd moments are never treated as probability moments.
All these certificates use exact arithmetic and are independent of
optimizer success or a sampled grid.

Let $\overline I_m,\overline O_m,\overline N_m$ denote the resulting
upper bounds, combined with the retained capped-mass bounds when
those are smaller. For even $m$ and $t\ge1$, the three positive
physical partition functions give the input bounds
\begin{align*}
 D&=(\overline I_m+2\overline O_m+3\overline N_m)^t,\\
 E_P&=(\overline I_m+2\overline O_m-\underline N_m)^t,\\
 E_C&=(\overline I_m-\underline O_m)^t,\\
 I_{\rm cap}&=(D+3E_P+8E_C)/12,\qquad
 J_{\rm cap}=(D+3E_P)/4.
\end{align*}
The chosen radicands are positive. This uses
$Z_m(tb_*,g)\le Z_m(b_*,g)^t$ for each twist, before applying
the sector inversion. No analogous time-power inequality is assumed
directly for a spatial sector.

The retained 120-site trial forces a unique dominant spatial
eigenvalue through the same capped residual comparison as before.
Its positivity follows from the signed trace at large odd lengths.
The 88-site trial supplies the physical-purity initialization.
The selected starting point is
\[
 m=28,\qquad t=41/20,\qquad n_0=44,\qquad b_0=2009/80.
\]
The exact seed and all intermediate fractions are recorded in
\path{independent_results/resolved_refinement.json}.

\subsection{A tighter conversion from purity to a gap}
\begin{lemma}\label{lem:tight-purity}
Let a physical Gibbs state at inverse temperature $b$ have purity
$T\ge K>1/2$. For any $E\ge e^{b\delta}>1$, the inequality
\[
 K>\frac{1+E^2}{(1+E)^2}
\]
implies a unique ground state and spectral gap greater than $\delta$.
\end{lemma}
\begin{proof}
The largest Gibbs weight $p_0$ is at least $T$, so $p_0>1/2$ and
the ground state is unique. Moreover
$T\le p_0^2+(1-p_0)^2$, since the remaining squared weights sum to
at most the square of their sum. The increasing branch of this
quadratic implies $p_0\ge(1+\sqrt{2K-1})/2$.
For the first excited weight $p_1$,
\[
 e^{-b\gap}=\frac{p_1}{p_0}
 \le\frac{1-p_0}{p_0}
 \le\frac{1-\sqrt{2K-1}}{1+\sqrt{2K-1}}.
\]
The displayed hypothesis gives a strict bound below $E^{-1}$,
which is at most $e^{-b\delta}$. The rational threshold is increasing
in $E>1$, so an upper Taylor enclosure for the exponential suffices.
\end{proof}
The two-level probability distribution saturates this conversion.
No assumption about the spin or degeneracy of the first excitation
is required. This elementary inequality improves finite-length
coverage; it does not, by itself, improve the asymptotic decay rate.

\subsection{Coverage and limits of the computation}
The sector recurrence and length interpolation in the next section
are retained, using even spatial reference lengths throughout.
The new certificate performs 304 named checks, with additional root-direction
and polynomial-positivity comparisons. It checks each length $24$ through
$65$ directly and then the following intervals:
\begin{center}
\begin{tabular}{rrr}
\toprule
First length & Last length & Recurrence level \\
\midrule
66 & 80 & 0 \\
81 & 246 & 0 \\
247 & 821 & 2 \\
822 & 2919 & 3 \\
2920 & 13012 & 5 \\
13013 & 33791 & 7 \\
\bottomrule
\end{tabular}
\end{center}
The tail starts at $L=33792$. At reference length $22528$ and inverse
temperature $64288/5$, the invariant holds with
$w_0=1.05443040418529\times10^{-811}$. Its update is $w_+=6w^2$.
The retained exponential comparisons give the strict gap bound on the
entire tail. The short-chain ladder is checked separately, including
$T(19,49/4)>.50776$, which now exceeds one half.
The ninth update requires outward products on a $10^{-1000}$ grid.
Root proposals use a smaller margin $10^{-450}$ and are still accepted
only after exact powered comparisons; the retained checkers keep
their original arithmetic settings. Exponential upper bounds use a rational
Taylor bound at an argument at most one, followed by upward-rounded
squaring to 40 significant digits. Every rounding direction is explicit.

The temperature exploration that motivated this revision is
non-rigorous. With the available moment orders, simply lowering the
temperature made the coarse two-sector initialization worse in the
screened examples. Resolving all three sectors was more effective.
Exploratory decay rates are neither uniform gap certificates nor
optimality barriers. The new theorem remains conditional on the
upstream transfer construction and has not received independent
external review or machine-formal verification.

\newpage
